\subsubsection{Distributed Representations}\label{sec:distributed-representations}

In the previous section we introduced the concept of a word embedding as a mapping from words to vectors in a continuous space. We can now make this idea more precise.

\highspace
Consider a vocabulary of $V$ words:
\begin{equation*}
    V = \{w_1, w_2, \ldots, w_{\left|V\right|}\}
\end{equation*}
In a one-hot representation (\autopageref{def:one-hot-encoding}), each word is associated with a unique position:
\begin{equation*}
    \text{Obama} = \left[0, \dots, 0, 1, 0, \dots, 0\right]
\end{equation*}
Only one component contains information. For this reason, \textbf{one-hot vectors} are often called \definition{Local Representations}: \hl{each word is identified by a single location in the vector.}

\highspace
In a \definition{Distributed Representation}, instead, \hl{a word is represented by a vector of real numbers}:
\begin{equation*}
    \text{Obama} = \left[0.12, -0.25, 0.81, \dots, 0.03\right]
\end{equation*}
Now \textbf{information is spread} across all dimensions. Each component contributes a small amount of information about the word. This is why the representation is called \textbf{distributed}. With local representations, semantic similarity is not captured, while with distributed representations, semantically similar words are mapped to nearby points in the vector space.

\highspace
\textcolor{Green3}{\faIcon{question-circle} \textbf{What does each position in the vector mean?}} In short, we usually \textbf{do not know} what each dimension means individually. And this is the fundamental difference between one-hot vectors and embeddings.
\begin{itemize}
    \item In a \textbf{one-hot vector} is easy to interpret because we know that the $i$-th position corresponds to the $i$-th word in the vocabulary. But it is not useful for learning.
    \item In a \textbf{distributed representation} is hard to interpret because the \hl{dimensions are learned automatically by the neural network.} The coordinates of the vector are not assigned by us, but are discovered automatically by the model during training. A dimension may partially encode gender, profession, nationality, formality, syntactic role, all mixed together.
\end{itemize}

\highspace
\begin{flushleft}
    \textcolor{Green3}{\faIcon{check-circle} \textbf{Fighting the Curse of Dimensionality}}
\end{flushleft}
\hl{Distributed Representation} is not only a different way of represeting words; it also \hl{addresses one of the major limitations} of traditional Natural Language Processing (NLP) representations: the \textbf{curse of dimensionality}. This goal is achieved through three mechanisms:
\begin{enumerate}
    \item \important{Compression}. A typical vocabulary may contain $\left|V\right| > 10^{6}$ distinct words. A one-hot representation therefore requires more than one million dimensions. A word embedding instead maps each word to a vector of dimension $100 < m < 500$. Thus $\mathbb{R}^{\left|V\right|} \to \mathbb{R}^{m}$ with $m \ll \left|V\right|$ mapping is a compression of the original representation.
    
    
    \item \important{Smoothing}. One-hot encodings are inherently discrete (i.e., they are either 0 or 1). There is \textbf{no notion of partial similarity}. Instead, embeddings transform words into points of a continuous space:
    \begin{equation*}
        C\left(w\right) \in \mathbb{R}^{m}
    \end{equation*}
    This process is called \definition{Smoothing} because \textbf{nearby points can represent similar concepts}.
    
    
    \item \important{Densification}. One-hot vectors are extremely sparse, and $\frac{1}{\left|V\right|}$ of the components are non-zero. For a vocabulary of one million words $\frac{1}{10^{6}} = 0.0001\%$, almost the entire vector contains no information. Instead, embeddings \textbf{transform} these \textbf{sparse vectors into dense feature vectors} where \hl{most components contribute information.}
\end{enumerate}
